\section{Current Grothendieck--Teichmüller scope: 20 November
2026}\label{current-grothendieckteichmuxfcller-scope-20-november-2026}

Research addendum \texttt{consequences\_criteria}; GPT-6 Astra Pro;
effort setting not exposed. This updates the literature status without
enlarging the adopted Alpöge input or claiming a kernel check of another
proof archive.

\subsection{CC-GT01. The 2026 rational claim has substantive primary
support}\label{cc-gt01.-the-2026-rational-claim-has-substantive-primary-support}

It would be outdated to label rational Deligne--Drinfeld freeness simply
open. OpenAI's 23 September 2026 preprint, Theorem 1.1, asserts freeness
of the rational graded solution Lie algebra on one generator in each odd
weight at least three, together with its weight-completed version. Its
mirrored introduction, the statement of its all-weight upper bound, and
parts of the argument were inspected; the entire 41-page original proof
and its Lean archive were not independently checked here. Plectis Atlas
itself is an index of other parties' reports, not our verification
method. The Valency page could not be retrieved directly.

Two subsequent primary mathematical texts supply stronger evidence than
an index entry:

\begin{longtable}[]{@{}
  >{\raggedright\arraybackslash}p{(\columnwidth - 4\tabcolsep) * \real{0.3333}}
  >{\raggedright\arraybackslash}p{(\columnwidth - 4\tabcolsep) * \real{0.3333}}
  >{\raggedright\arraybackslash}p{(\columnwidth - 4\tabcolsep) * \real{0.3333}}@{}}
\toprule\noalign{}
\begin{minipage}[b]{\linewidth}\raggedright
Source
\end{minipage} & \begin{minipage}[b]{\linewidth}\raggedright
Exact locator inspected
\end{minipage} & \begin{minipage}[b]{\linewidth}\raggedright
Mathematical contribution
\end{minipage} \\
\midrule\noalign{}
\endhead
\bottomrule\noalign{}
\endlastfoot
Willwacher (2026b), arXiv:2610.15592v4, 28 October & Theorem 1.1, p.~2;
Theorem 2.2, p.~3, proved in §§3.1--3.3, pp.~4--6 & The required upper
bound \(\dim_{\mathbb Q}\mathfrak{grt}_1(\mathbb Q)_n\le d_n\), where
\(d_n\) is the dimension of the free Lie algebra on odd generators \\
Hirose (2026), arXiv:2610.23719v1, 16 October & Theorem 4 and §1.2, p.~2
& An integral Hoffman spanning theorem modulo confluence relations,
combined there with Brown and Furusho to identify the motivic group with
\(\mathrm{GRT}\) \\
\end{longtable}

For the first paper, the full chain proving the upper bound was read:
integral rank comparison, the characteristic-two projection, its
explicit reconstruction, the image restriction, and the free-Lie
dimension calculation. The key derivation intertwining and square-zero
identities were checked algebraically. This is a targeted reading of an
existing proof, not an independent reproof or a formal verification.
Theorem 1.3 and §4 give further characteristic-two structure, but they
are unnecessary when Brown's injection is retained. Hirose's theorem and
stated motivic consequence were inspected; his full technical proof was
not checked in this run.

Neither proof uses RH or Dirichlet GRH. They are external 2026
mathematical results, not achievements attributable to the adopted
analytic input.

\subsection{CC-GT02. What Drinfeld and Brown had already
supplied}\label{cc-gt02.-what-drinfeld-and-brown-had-already-supplied}

Drinfeld (1991), Proposition 6.3, p.~859, gives a homogeneous element of
\(\mathfrak{grt}_1\) in every odd degree \(n\ge3\) with nonzero
depth-one term. The unnumbered ``Questions'' immediately following ask
about precisely one generator in each odd degree and freeness. The
primary proposition, proof and questions were visually inspected on the
rendered original scan; its existence assertion should not be conflated
with freeness or exhaustion.

Brown (2012), Theorem 1.1, p.~950, gives the motivic Hoffman basis; the
paragraph immediately following explains why it proves Conjecture 1 on
p.~949, namely faithfulness of the motivic action. The structure of the
motivic prounipotent Lie algebra is recalled on p.~950. An exact
numbered modern formulation of its free-subalgebra consequence is
Willwacher (2026a), Theorem 1.1, p.~1. His Theorem 1.2, also p.~1,
extends the assertion to any homogeneous odd-degree family with nonzero
depth-one coefficients. These are injectivity statements. The upper
bound in CC-GT01 supplies the distinct exhaustion step.

\subsection{CC-GT03. Combining injection and upper bound, including
completion}\label{cc-gt03.-combining-injection-and-upper-bound-including-completion}

\textbf{Statement.} Let \[
 \mathfrak f=\operatorname{Lie}_{\mathbb Q}\langle e_3,e_5,e_7,\ldots\rangle,
 \qquad \deg e_j=j.
\] Write
\(\mathfrak g^{\mathrm{gr}}=\bigoplus_n\mathfrak{grt}_{1,n}(\mathbb Q)\)
for the graded algebra of polynomial solutions, and
\(\widehat{\mathfrak g}=\prod_n\mathfrak{grt}_{1,n}(\mathbb Q)\) for its
weight completion; some sources use \(\mathfrak{grt}_1\) for the latter.
Take the weight-preserving injection
\(\iota:\mathfrak f\hookrightarrow\mathfrak g^{\mathrm{gr}}\) supplied
by Brown's result and the upper bound of Willwacher (2026b), Theorem
1.1. Then \(\iota\) is an isomorphism, and its extension
\(\prod_n\mathfrak f_n\to\widehat{\mathfrak g}\) is an isomorphism of
complete Lie algebras.

\textbf{Proof.} In every weight \(n\), the space \(\mathfrak f_n\) is
finite-dimensional: only finitely many generators can occur, and every
generator has degree at least three. Injectivity gives \[
 \dim\mathfrak f_n\le\dim\mathfrak{grt}_{1,n}.
\] The cited upper bound reverses this inequality. Therefore the
injective linear map \(\iota_n\) is bijective in every weight. Their
direct sum is a bijective graded Lie homomorphism, hence a Lie algebra
isomorphism. Their product is a continuous bijection for the product,
equivalently weight-filtration, topologies; its inverse is the product
of the degreewise inverses and is continuous. A bracket component of
fixed weight is a finite sum of brackets of lower-weight components, so
the product map and its inverse preserve the completed bracket. This
proves the completed assertion. \(\square\)

This deduction identifies exactly which result completed the rational
theorem. It supplies no theorem about the full profinite absolute Galois
group.

\subsection{CC-GT04. The full profinite question is a different
assertion}\label{cc-gt04.-the-full-profinite-question-is-a-different-assertion}

Fix Drinfeld's original profinite group \(\widehat{\mathrm{GT}}\), with
its commutator, invertibility, symmetry, hexagon and pentagon
conditions. Bortnovskyi, Dolgushev, Holikov, and Pashkovskyi (2026
revision), §1.3, equations (1.5)--(1.8), pp.~4--5, gives the fixed Ihara
embedding and cyclotomic compatibility. In that paper's coordinates, \[
 \operatorname{Ih}(g)=\bigl((\chi_{\mathrm{cyc}}(g)-1)/2,f_g\bigr),
 \qquad
 x\mapsto x^{\chi_{\mathrm{cyc}}(g)},\quad
 y\mapsto f_g^{-1}y^{\chi_{\mathrm{cyc}}(g)}f_g.
\] Its cyclotomic projection is \((m,f)\mapsto2m+1\), so the composite
equals \(\chi_{\mathrm{cyc}}\). The paper's introduction distinguishes
surjectivity of this whole embedding from its proved surjectivity onto
particular quotients; its Theorems 5.3 and 5.7 concern such quotient
results. Rational Lie-algebra freeness does not establish full profinite
surjectivity.

The safe status statement for this report is precise: \textbf{the full
profinite identification is not settled by the adopted RH/Dirichlet GRH
or by the rational Deligne--Drinfeld theorem; the inspected 2026 primary
profinite work treats full surjectivity as unresolved.} No exhaustive
claim about all unpublished assertions is needed. The gentle,
pro-\(\ell\), prounipotent, and higher-genus variants must not be
substituted for the fixed original profinite embedding.

\subsection{CC-GT05. The exact positive role of the cyclotomic
quotient}\label{cc-gt05.-the-exact-positive-role-of-the-cyclotomic-quotient}

The cyclotomic character and Kronecker--Weber identify finite abelian
characters of \(G_{\mathbb Q}\) with Dirichlet characters. Consequently
the adopted GRH controls the analytic factors of finite abelian Artin
representations and their Tate twists, including the Borger realizations
of CC-F01 and CC-F01b. Those results also apply to arithmetic
representations known to factor through the cyclotomic projection in the
preceding diagram.

This analytic statement concerns those representations and their
\(L\)-functions. It does not determine the kernel of the cyclotomic
projection, the associator component \(f_g\), or surjectivity of
\(\operatorname{Ih}\). Abstractly, surjectivity after projection cannot
establish surjectivity before projection: for any profinite abelian
group \(A\), the map \(A\to A\times C_2\), \(a\mapsto(a,1)\), has
surjective composite with the projection to \(A\) and is a proper
embedding. This elementary observation is a counterexample to that
general inference, not a counterexample to any conjecture about the
actual group \(\widehat{\mathrm{GT}}\).

\subsection{Exact bibliography and access
record}\label{exact-bibliography-and-access-record}

\begin{enumerate}
\def\labelenumi{\arabic{enumi}.}
\tightlist
\item
  \textbf{Brown, F. (2012).} Mixed Tate motives over \(\mathbb Z\).
  \emph{Annals of Mathematics, 175}(2), 949--976.
  https://doi.org/10.4007/annals.2012.175.2.10 . Primary PDF:
  https://annals.math.princeton.edu/wp-content/uploads/annals-v175-n2-p10-p.pdf
  . Theorem 1.1 and its stated consequence directly inspected.
\item
  \textbf{Drinfeld, V. G. (1991).} On quasitriangular quasi-Hopf
  algebras and on a group that is closely connected with
  \(\operatorname{Gal}(\overline{\mathbb Q}/\mathbb Q)\).
  \emph{Leningrad Mathematical Journal, 2}(4), 829--860. Original
  Russian publication: \emph{Algebra i Analiz, 2}(4), 149--181 (1990).
  Scan: https://www.lamfa.u-picardie.fr/marin/docs/drinfeld.pdf .
  Proposition 6.3 and the following unnumbered Questions, p.~859,
  visually inspected on the rendered original page.
\item
  \textbf{Willwacher, T. (2026a).} \emph{On free generators in the
  Grothendieck--Teichmüller Lie algebra} (arXiv:2607.18793v2, 6 August
  2026). https://arxiv.org/pdf/2607.18793v2 . Theorems 1.1--1.2, p.~1,
  directly inspected.
\item
  \textbf{Willwacher, T. (2026b).} \emph{The Deligne--Drinfeld
  conjecture after OpenAI} (arXiv:2610.15592v4, 28 October 2026).
  https://arxiv.org/pdf/2610.15592v4 . Theorem 1.1 and its entire proof
  chain in §§1--3 read; key algebraic identities checked. Theorem 1.3
  and §4 are additional, not required inputs to CC-GT03.
\item
  \textbf{Hirose, M. (2026).} \emph{Confluence relations, Hoffman basis
  over \(\mathbb Z_{(2)}\), and the Deligne--Drinfeld conjecture}
  (arXiv:2610.23719v1, 16 October 2026).
  https://arxiv.org/pdf/2610.23719v1 . Theorem 4 and §1.2, p.~2,
  directly inspected; full technical proof not independently checked
  here.
\item
  \textbf{Bortnovskyi, I., Dolgushev, V. A., Holikov, B., \&
  Pashkovskyi, V. (2024/2026).} \emph{Accessing non-abelian quotients of
  the Grothendieck--Teichmueller group via elementary tools}
  (arXiv:2405.11725v3, 6 August 2026).
  https://arxiv.org/pdf/2405.11725v3 . Introduction, §1.3, equations
  (1.5)--(1.8), pp.~4--5, and the described quotient scope directly
  inspected.
\item
  \textbf{OpenAI. (2026, September 23).} \emph{The Deligne--Drinfeld
  conjecture}. OpenAI Math Release preprint.
  https://github.com/openai/math/blob/fd4aeeb2ee4fc729c18d98444fed42fd0529eeeb/preprints/The-Deligne-Drinfeld-conjecture-September-23-2026/paper.pdf
  . Inspected mirror:
  https://wcook04.github.io/plectis-atlas/papers/The-Deligne-Drinfeld-conjecture-September-23-2026.html
  . Theorem 1.1 and dimension-bound route inspected; the complete
  original proof and Lean archive were not checked in this run.
\end{enumerate}
